\section{Introduction}\label{sec:intro}

The Internet of Things (IoT) is transforming residential environments
into complex cyber-physical systems.  A modern smart home may contain
dozens of heterogeneous devices---sensors, actuators, voice assistants,
and hub controllers---whose interactions with human occupants produce
rich, privacy-sensitive data streams.  Researchers studying smart home
security~\cite{iot-security-survey}, energy optimization~\cite{iot-energy-survey}, and activity
recognition~\cite{har-survey} all share a common bottleneck: the need for
\emph{realistic, large-scale, and reproducible} testbeds.

Physical smart home testbeds~\cite{iot-testbed-survey} are expensive to construct,
difficult to instrument, and inherently limited to a single floor plan
and occupant population.  Software-only simulators such as
OpenHAB-Sim~\cite{openhab-sim} and IoTSim~\cite{iotsim} improve reproducibility but
sacrifice fidelity: they model devices as stateless message endpoints,
ignore firmware-level behavior, and generate occupant activity through
simple Markov chains that fail to capture the contextual richness of
real daily routines.

We present \sys{} (\textbf{V}irtual \textbf{E}nvironment for
\textbf{S}mart-home \textbf{P}latform \textbf{E}valuation and
\textbf{R}esearch), a simulation platform that addresses these
limitations through three tightly integrated innovations:

\begin{enumerate}[leftmargin=*]
  \item \textbf{Firmware-in-the-loop emulation.}
    \sys{} runs real ARM device firmware inside QEMU containers
    managed by Docker, enabling bit-accurate execution of sensor
    sampling loops, interrupt handlers, and communication stacks
    (UART-over-TCP, MQTT).  Each virtual device exposes a network
    interface matching its physical counterpart's protocol.

  \item \textbf{LLM-driven activity generation.}
    Rather than relying on hand-crafted Markov models, \sys{} uses
    large language models (via LMStudio) to generate daily activity
    schedules conditioned on detailed occupant personas (age,
    occupation, mobility constraints, lifestyle preferences).  This
    produces statistically diverse yet contextually coherent
    behavioral patterns.

  \item \textbf{Bi-directional cloud synchronization.}
    \sys{} implements the Samsung SmartThings Schema Connector
    protocol, allowing simulated devices to appear as first-class
    citizens in a real SmartThings account.  State changes flow
    in both directions, enabling Sim2Real validation of automations,
    routines, and third-party integrations.
\end{enumerate}

These components are unified within Meta's Habitat~3.0 embodied AI
framework, which provides physics-accurate 3D scene rendering,
humanoid navigation, and object interaction.  A central event bus
coordinates real-time data flow among all subsystems at
sub-millisecond latency.

\paragraph{Contributions.}
\begin{itemize}[leftmargin=*]
  \item The design and implementation of \sys{}, a smart home
    simulation platform that, to our knowledge, is the first to
    combine embodied 3D simulation, firmware-in-the-loop emulation,
    and cloud-connected IoT in a single reproducible framework.
  \item An LLM-based activity generation pipeline with a library of
    10~diverse personas, producing schedules whose distributions
    show moderate alignment with real-world datasets (CASAS, ARAS)
    with a mean JS divergence of~0.218 across five reference homes,
    outperforming a learned Markov baseline on diverse multi-occupant
    environments.
  \item A comprehensive evaluation across five research questions,
    demonstrating fidelity (with baseline comparison), scalability,
    latency, Sim2Real transfer potential (0.813~accuracy on ARAS-A
    via enriched 119-feature representations), and security
    assessment capability (validated against 24~real-world CVEs
    covering 7 of 10 OWASP IoT Top~10 categories).
  \item A large-scale autonomous evaluation across 28~HSSD scenes
    over 7~simulated days each (88.0~hours wall-clock), validating
    end-to-end reliability with 94.9\% navigation success,
    47{,}207~firmware-triggered cloud updates, and
    4{,}307~LLM-generated tasks.
  \item A security testing framework comprising 36~unique attacks
    across five suites (firmware, network, phantom-delay, malicious
    SmartApp, ESP32 buffer overflow) with automated CVSS~3.1 scoring
    and MITRE ATT\&CK mapping, yielding the first IoT security
    evaluation integrated within an embodied simulation platform.
    Packet-level traffic analysis---validated by \texttt{tshark}
    pcap captures (154{,}151~TCP segments, 12.1\,MB global capture,
    736~individual pcap files)---confirms that all attack traffic
    traverses the real Docker bridge network and is independently
    verifiable with Wireshark.
  \item An open-source release of the platform, evaluation framework,
    and all experiment configurations for full reproducibility.
\end{itemize}
